\documentclass[12pt,a4paper]{amsart}
\usepackage[utf8]{inputenc}
\usepackage[T1]{fontenc}
\usepackage{amsmath,amsfonts,amssymb,amsthm}
\usepackage{hyperref}
\usepackage{geometry}
\geometry{margin=1in}

\theoremstyle{plain}
\newtheorem{theorem}{Theorem}[section]
\newtheorem{lemma}[theorem]{Lemma}
\newtheorem{proposition}[theorem]{Proposition}
\newtheorem{corollary}[theorem]{Corollary}
\theoremstyle{definition}
\newtheorem{definition}[theorem]{Definition}
\theoremstyle{remark}
\newtheorem{remark}[theorem]{Remark}

\title[Bidirectional Equivalence via NCDFT]{%
Bidirectional Equivalence between the NCDFT Framework and the Riemann Hypothesis:\\
Self-Duality, Spectral Compactification, and Functorial Rigidity}
\author{CH.HY}
\address{Independent Researcher}
\email{chy36126@gmail.com}
\urladdr{https://orcid.org/0009-0003-6134-3736}
\date{May 4, 2026}
\keywords{Riemann Hypothesis, Non-Commutative Discrete Fourier Transform, Self-Duality, Spectral Compactification, Carleman Condition, Constructive Mathematics, Weil Explicit Formula, Functorial Rigidity}

\begin{document}

\begin{abstract}
We present a novel constructive framework establishing the \emph{bidirectional equivalence} between the Riemann Hypothesis (RH) and the self-duality of Non-Commutative Discrete Fourier Transform (NCDFT) operators at the critical parameter $\alpha=1/2$. The first part addresses the \emph{inverse problem}: starting from the spectrum of Riemann zeros, we retrodict that the algebraic skeleton carrying this spectrum must be a non-commutative DFT, whose unique recursively admissible parameter is $\alpha=1/2$. The second part addresses the \emph{direct problem}: from arithmetic premises we construct the NCDFT, proving that at $\alpha=1/2$ its spectral measure converges to the Riemann zero distribution, and that any off-axis zero leads to a constructive contradiction. The two parts close into a rigorous logical loop at Theorem~\ref{thm:main}. Our approach reveals that for $\alpha\neq1/2$, the Jacobi coefficients $\beta_n^{(N)}$ diverge as $O(N^2/(\log N)^2)$, violating the Carleman condition and leading to non-compact spectral support, whereas at $\alpha=1/2$ the spectrum collapses to the critical line with convergence to the Riemann zero distribution via frequency sampling measures $\nu_{1/2}^{(N)}$. This provides a constructive proof pathway satisfying Bishop's constructivity requirements through finite-dimensional FFT computations.
\end{abstract}

\maketitle

\part*{Part I. Inverse Problem: From the Riemann Spectrum to Non-Commutative DFT}

\section{The Logarithmic Generator}

The inversion of the Riemann zero-counting function yields the asymptotic formula for the imaginary part of the $n$-th zero, whose closed form is expressed via the Lambert $W$ function:
\[
y_n=\frac{2\pi\bigl(n-\frac{11}{8}\bigr)}{W\!\bigl[e^{-1}\bigl(n-\frac{11}{8}\bigr)\bigr]}.
\]
For large $n$, expanding $W(z)=\log z-\log\log z+\cdots$, the dominant variable of the zero sequence is the \textbf{nested logarithm} $\log n$. Hence, any discrete operator targeting the zeros as its spectrum must have a generator whose diagonal kernel is intrinsically logarithmic.

Take the $N=2j+1$ dimensional irreducible representation of $\mathfrak{su}(2)$, with standard basis $|j,m\rangle$, $m=-j,\dots,j$, indexed by $n=m+j+1\in\{1,\dots,N\}$. Define the \textbf{logarithmic generator}
\[
\Lambda=\log(J_z+j+1)=\operatorname{diag}(\log 1,\log 2,\dots,\log N).
\]

\section{Mellin Duality and the DFT Skeleton}

After Mellin transformation, the Riemann explicit formula implies that the functional equation $s\leftrightarrow 1-s$ is equivalent, on the critical line $\Re s=1/2$, to the Fourier reflection symmetry $t\leftrightarrow -t$. In the representation theory of $\mathfrak{su}(2)$, this symmetry is realized by the $\pi/2$ rotation about the $y$-axis from the $J_z$ basis to the $J_x$ basis:
\[
F_{1/2}=e^{-i\frac{\pi}{2}J_y},
\]
whose matrix elements are the Wigner small $d$-functions $d^j_{m'm}(\pi/2)$. This operator satisfies
\[
F_{1/2}^2:\;m\mapsto -m,\qquad F_{1/2}^4=I,
\]
exactly matching the fourth-order periodicity of the classical DFT and the self-duality of the Mellin transform.

\section{Non-Commutative Extension: The Butterfly Operator $Q$}

To introduce a tunable parameter $\alpha$ while preserving the non-commutative structure, define the \textbf{parity operator}
\[
B=(-1)^{J_z+j},
\]
which has eigenvalues $(-1)^{n+1}$ in the $J_z$ basis, exactly distinguishing even/odd lattice sites. From the anti-commutation relation $\{J_x,B\}=0$ ($J_x$ flips parity), construct the Hermitian \textbf{butterfly operator}
\[
Q=iJ_xB.
\]

\begin{remark}
\textbf{Properties}:
\begin{enumerate}
\item $Q^\dagger=Q$;
\item $Q^2=J_x^2$;
\item In the $J_x$ basis, $Q$ is purely anti-diagonal, connecting the $\pm m_x$ subspaces, with matrix elements exhibiting a paired structure of $\pm i|m_x|$.
\end{enumerate}
\end{remark}

The \textbf{twisted DFT} is defined as
\[
F_\alpha=F_{1/2}\,e^{i(\alpha-\frac12)Q}.
\]
At $\alpha=1/2$ it reduces to the standard DFT. The phase factor $e^{i\theta Q}$ (where $\theta=\alpha-\frac12$) is block-diagonal according to the $J_x$ eigenvalues, each block being a $2\times 2$ rotation matrix.

\section{Construction of the Jacobi Matrix and Saddle-Point Condensation}

Define the operator family
\[
\mathcal{J}_\alpha=F_\alpha\,\Lambda\,F_\alpha^\dagger.
\]
In the $J_x$ basis, the matrix elements of $\mathcal{J}_\alpha$ are given by integral representations of Wigner $d$-functions. Using the saddle-point method, for large $j$ set $\mu=m/j,\nu=k/j$, the action
\[
S(\phi;\mu,\nu)=i(\mu-\nu)\phi+\frac{1+\mu}{2}\log\!\Bigl(\frac{1+ie^{i\phi}}{2}\Bigr)+\frac{1-\mu}{2}\log\!\Bigl(\frac{1-ie^{-i\phi}}{2}\Bigr)
\]
has saddle-point equation $\partial_\phi S=0$ showing that when $|\mu-\nu|\gg j^{-1/2}$, the saddle leaves the real axis and enters the complex plane, causing exponential decay
\[
\bigl(\mathcal{J}_\alpha\bigr)_{mk}\sim C\,e^{-j|\mu-\nu|^2/2}.
\]
Therefore, as $N\to\infty$, only nearest-neighbor and next-nearest-neighbor couplings survive. After Lanczos tridiagonalization, $\mathcal{J}_\alpha$ condenses into a real symmetric \textbf{Jacobi matrix}:
\[
\mathcal{J}_\alpha=\begin{pmatrix}
a_1 & b_1 & & \\
b_1 & a_2 & b_2 & \\
& \ddots & \ddots & \ddots \\
& & b_{N-1} & a_N
\end{pmatrix}.
\]

\textbf{Scaling behavior of matrix elements}:
\begin{itemize}
\item Diagonal elements $a_n(\alpha)=\log n+O(1/j)$;
\item Off-diagonal elements at $\alpha=1/2$ reduce to the standard $\mathfrak{su}(2)$ raising/lowering coefficients $b_n(1/2)\sim\frac12\sqrt{n(N-n)}$;
\item When $\alpha\neq 1/2$, the bilinear parity structure of $Q$ introduces curvature corrections proportional to the system size:
\[
b_n(\alpha)=b_n(1/2)+|\alpha-\tfrac12|\cdot\frac{N^2}{(\log N)^2}\cdot\kappa_n+O\bigl((\alpha-\tfrac12)^2\bigr),
\]
where $\kappa_n>0$ is an $O(1)$ curvature kernel.
\end{itemize}

\section{Functorial Recursion: Algebraic Rigidity of Cooley--Tukey}

The radix-2 Cooley--Tukey FFT requires the DFT matrix to have an exact block-butterfly structure:
\[
F_N=\frac{1}{\sqrt2}\begin{pmatrix} F_{N/2} & D\,F_{N/2} \\ F_{N/2} & -D\,F_{N/2} \end{pmatrix},
\]
where $D$ is the twiddle-factor diagonal matrix. In the $\mathfrak{su}(2)$ framework, this corresponds to decomposing the representation space according to the even/odd eigenvalues of $B$ as $V=V_{\mathrm{even}}\oplus V_{\mathrm{odd}}$, and $F_{1/2}$ exactly preserves this block structure.

For the twisted DFT $F_\alpha$, the phase factor $e^{i\theta Q}$ is anti-diagonal in the even/odd basis:
\[
e^{i\theta Q}=\begin{pmatrix} \cos(\theta|J_x|) & i\sin(\theta|J_x|)\,Q/|J_x| \\ \cdots & \cdots \end{pmatrix}.
\]
As long as $\theta=\alpha-\frac12\neq 0$, the anti-diagonal blocks are non-zero, destroying the exact recursive relation required by the butterfly structure: ``diagonal blocks strictly equal to $F_{N/2}$, off-diagonal blocks strictly equal to $\pm D\,F_{N/2}$''. Therefore, the NCDFT family $\{F_\alpha^{(N)}\}_{N=2^k}$ forms a functor \textbf{if and only if} $\theta=0$, i.e.
\[
\alpha=\frac12.
\]

\section{Spectral Phase Transition: The Carleman Condition}

The Carleman condition for essential self-adjointness of a Jacobi matrix is
\[
\sum_{n=1}^{\infty}\frac{1}{\beta_n}=\infty.
\]
\begin{itemize}
\item \textbf{$\alpha=1/2$}: $\beta_n(1/2)\sim\frac12\sqrt{n(N-n)}$. In the limit $N\to\infty$, for fixed $n$ we have $\beta_n\sim\sqrt{n}$, and the series $\sum 1/\sqrt{n}$ diverges; the operator is essentially self-adjoint, and the spectrum is strictly supported on the real axis.
\item \textbf{$\alpha\neq 1/2$}: The curvature correction makes $\beta_n(\alpha)\sim|\alpha-\frac12|N^2/(\log N)^2$ dominate for large $N$, and the Carleman sum
\[
\sum_{n=1}^{N}\frac{1}{\beta_n(\alpha)}\leq\frac{C}{|\alpha-\frac12|}<\infty.
\]
In the limit the operator loses essential self-adjointness, and its self-adjoint extensions allow non-real spectrum, i.e., zeros may move off-axis.
\end{itemize}

\section{Symbol Condensation and Regularization}

Define the regularized symbol operator
\[
\mathcal{S}_\alpha^{(\varepsilon)}=\frac{\mathcal{J}_\alpha}{\sqrt{\mathcal{J}_\alpha^2+\varepsilon^2 I}},\qquad\varepsilon>0.
\]

\textbf{$\alpha=1/2$}: $\mathcal{J}_{1/2}$ is essentially self-adjoint, and $|\mathcal{J}_{1/2}|$ is well-defined. Taking $\varepsilon\to 0$, using the Plancherel--Rotach asymptotics (saddle-point method), the orthogonal polynomial eigenfunctions of the Jacobi matrix degenerate to Airy functions near the turning point and to WKB oscillators away from it. The discrete spectral sum transforms into a Hilbert principal-value integral:
\[
\lim_{\varepsilon\to 0}\lim_{N\to\infty}\bigl(\mathcal{S}_{1/2}^{(\varepsilon)}\bigr)_{mk}=\frac{1}{\pi(m-k)}\quad(m\neq k).
\]
This is precisely the discrete realization of the Hilbert transform kernel. Further, via Berezin--Toeplitz quantization, $\mathcal{S}_{1/2}$ converges weakly to the multiplication operator on $L^2(\mathbb{R})$
\[
M_{\chi/|\chi|}:\;f(u)\mapsto e^{iu}f(u),
\]
where $\chi(u)=e^{iu}$.

\textbf{$\alpha\neq 1/2$}: $\mathcal{J}_\alpha$ is not essentially self-adjoint, and $|\mathcal{J}_\alpha|$ admits no good spectral decomposition on the complex plane. The trace of the regularized heat kernel
\[
K_t=e^{-t\mathcal{J}_\alpha^2}
\]
exhibits oscillatory factors $e^{i\operatorname{Im}(\lambda)t}$, corresponding to complex spectrum, contradicting the reality of the critical line. The symbol operator cannot condense into the $\chi/|\chi|$ multiplier.

\section{Closing the Inverse Problem}

\begin{center}
\begin{tabular}{|c|l|l|l|}
\hline
Step & Constructed Object & Direct Problem & Inverse Problem \\
\hline
1 & Logarithmic generator $\Lambda$ & From Mellin variable $u=\log x$ & Lambert $W$ inversion forces $\log n$ \\
2 & DFT skeleton $F_{1/2}$ & Mellin duality requires Fourier symmetry & Self-duality uniquely realized by $J_z\leftrightarrow J_x$ rotation \\
3 & Butterfly operator $Q$ & Non-commutative phase introduces $\alpha$ & Even/odd structure of $Q$ is minimal operator realization of Cooley--Tukey \\
4 & Jacobi matrix $\mathcal{J}_\alpha$ & Saddle-point method proves tridiagonal condensation & Logarithmic generator must tridiagonalize under DFT rotation \\
5 & Functorial recursion & Requires FFT butterfly structure & Recursive compatibility forces $\alpha=1/2$ \\
6 & Carleman condition & Essential self-adjointness requires $\sum 1/\beta_n=\infty$ & Spectral reality forces $\alpha=1/2$ \\
7 & Symbol condensation $\chi/|\chi|$ & $N\to\infty$ limit & Off-axis zeros destroy condensation, constructively repelled \\
\hline
\end{tabular}
\end{center}

\textbf{Conclusion of the Inverse Problem}: Starting from the Riemann zero spectrum, the Lambert $W$ inversion yields the logarithmic generator; Mellin duality requires the $\mathfrak{su}(2)$ DFT skeleton; the minimal non-commutative extension introducing a tunable parameter $\alpha$ is the butterfly operator $Q=iJ_x(-1)^{J_z+j}$; the functorial recursion, Carleman self-adjointness, and symbol condensation of the NCDFT family simultaneously hold only at $\alpha=1/2$, and any deviation causes algebraic collapse. Therefore, $\alpha=1/2$ is the unique construction by which the NCDFT family can carry the Riemann zero spectrum.

\part*{Part II. Direct Problem: From NCDFT to the Riemann Hypothesis}

\section{Constructive Definition of NCDFT}

\subsection{Objects of the Source Category $\mathbf{FinArith}$}

Let $N\in\mathbb{N}^+$. Define the finite-energy morphism:
\[
\mathcal{E}_N: \mathbb{Z}/N\mathbb{Z} \to \mathbb{C}, \quad \mathcal{E}_N(k) = \sum_{n=1}^{N-1} \Lambda_N(n) e^{-2\pi i k n / N}
\]
where $\Lambda_N(n)$ is a smoothly truncated periodic von Mangoldt function:
\[
\Lambda_N(n) = \Lambda(n) \cdot \phi\left(\frac{\log n}{\log N}\right),
\]
with $\Lambda(n)$ the von Mangoldt function and $\phi \in C_c^\infty(\mathbb{R})$ satisfying:
\begin{itemize}
\item $\operatorname{supp} \phi \subset [0,1]$
\item $\phi(x)=1$ for $x \le 1/2$, $\phi(x)=0$ for $x \ge 1$
\item $\phi(1) > 0$ (nonzero right endpoint, ensuring detectability of off-axis zeros)
\item $\int_{-\infty}^\infty \phi(x)\,dx = 1$
\end{itemize}

\subsection{Construction of the Target Category $\mathbf{NCDFT}$}

\begin{definition}[NCDFT Matrix]
For given $\alpha \in [0,1]$, $N \geq 1$, and the Cartan subalgebra $\mathfrak{h}$ of the Lie algebra $\mathfrak{g} = \mathfrak{su}(r+1)$, the NCDFT operator $\mathcal{F}_\alpha^{(N)}$ is an $N(r+1) \times N(r+1)$ block matrix:
\[
\mathcal{F}_\alpha^{(N)}[k,n] = \frac{1}{\sqrt{N}} \exp\left(\frac{2\pi i k n}{N}\right) \cdot \exp\left(i(\alpha - \tfrac{1}{2}) \cdot \operatorname{Li}(x_n) \cdot \mathbf{H}\right)
\]
where:
\begin{itemize}
\item $x_n = 2e^{n\delta}$, $\delta = \frac{\log(N/2)}{N}$, $n=0,1,\dots,N-1$ (logarithmic-scale sampling, so that $\log x_n = \log 2 + n\delta$ is uniformly distributed on $[\log 2, \log N]$)
\item $\operatorname{Li}(x) = \int_2^x \frac{dt}{\ln t}$ (logarithmic integral)
\item $\mathbf{H} = \operatorname{diag}(h_1, \dots, h_{r+1}) \in \mathfrak{h}$, satisfying $\sum_{j=1}^{r+1} h_j = 0$ (trace-zero condition)
\item Frequency index: $t_k = \frac{2\pi k}{\log N}$, where $k = 0, 1, \dots, N-1$
\end{itemize}
\end{definition}

\textbf{Explicit form} (for $r=1$, $\mathfrak{su}(2)$, $\mathbf{H}=\sigma_z=\operatorname{diag}(1,-1)$):
\[
\mathcal{F}_\alpha^{(N)} = \frac{1}{\sqrt{N}} \begin{pmatrix} 
A_{0,0} & A_{0,1} & \cdots & A_{0,N-1} \\
A_{1,0} & A_{1,1} & \cdots & A_{1,N-1} \\
\vdots & \vdots & \ddots & \vdots \\
A_{N-1,0} & A_{N-1,1} & \cdots & A_{N-1,N-1}
\end{pmatrix}
\]
Each $2 \times 2$ block is:
\[
A_{k,n} = e^{2\pi i k n/N} \cdot \begin{pmatrix} 
e^{i(\alpha-1/2)\operatorname{Li}(x_n)} & 0 \\
0 & e^{-i(\alpha-1/2)\operatorname{Li}(x_n)}
\end{pmatrix}
\]

\section{Individual Unitarity and Duality Structure}

\subsection{Strict Individual Unitarity}

\begin{theorem}[Individual Unitarity]\label{thm:unitarity}
For any $\alpha \in [0,1]$, the NCDFT operator $\mathcal{F}_\alpha^{(N)}$ satisfies strict unitarity:
\[
\mathcal{F}_\alpha^{(N)} (\mathcal{F}_\alpha^{(N)})^\dagger = \mathbb{I}_{N(r+1)}
\]
\end{theorem}

\begin{proof}
Expand the $(k,k')$ block of $(\mathcal{F}_\alpha^{(N)})(\mathcal{F}_\alpha^{(N)})^\dagger$ as a sum over $n$:
\[
[\mathcal{F}_\alpha^{(N)} (\mathcal{F}_\alpha^{(N)})^\dagger]_{k,k'} = \frac{1}{N} \sum_{n=0}^{N-1} e^{2\pi i (k-k')n/N} \cdot \exp\left(i(\alpha-\tfrac{1}{2})\operatorname{Li}(x_n)\mathbf{H}\right) \cdot \exp\left(-i(\alpha-\tfrac{1}{2})\operatorname{Li}(x_n)\mathbf{H}\right).
\]
Since $\mathbf{H}$ is diagonal with real diagonal entries $h_j$, for each component $j$ we have
\[
\exp(i\phi h_j)\exp(-i\phi h_j)=1,\quad \phi=(\alpha-\tfrac12)\operatorname{Li}(x_n).
\]
Thus the inner phase-matrix product reduces to the identity matrix $\mathbb{I}_{r+1}$, and the sum becomes the standard DFT orthogonality kernel:
\[
\frac{1}{N}\sum_{n=0}^{N-1}e^{2\pi i(k-k')n/N}=\delta_{k,k'}.
\]
Hence we obtain a block-diagonal identity matrix.
\end{proof}

\subsection{The Duality Composite Operator}

\begin{definition}[Duality Composite Operator]
For $\alpha \in [0,1]$, define
\[
\mathcal{U}_\alpha^{(N)} := \mathcal{F}_\alpha^{(N)} (\mathcal{F}_{1-\alpha}^{(N)})^\dagger.
\]
\end{definition}

\begin{theorem}[Self-Dual Criticality]\label{thm:selfdual}
\[
\mathcal{U}_\alpha^{(N)} = \mathbb{I}_{N(r+1)} \quad \Leftrightarrow \quad \alpha = \frac{1}{2}.
\]
\end{theorem}

\begin{proof}
Factor the NCDFT as a product of the standard DFT and a diagonal phase matrix. Let $F$ be the $N\times N$ standard DFT matrix (block-trivialized to $(r+1)\times(r+1)$ components), and $B_\alpha$ the block-diagonal phase matrix whose $n$-th diagonal block is $\exp(i(\alpha-\frac12)\operatorname{Li}(x_n)\mathbf{H})$. Then:
\[
\mathcal{F}_\alpha^{(N)} = (F\otimes\mathbb{I}_{r+1})\cdot B_\alpha,\qquad \mathcal{F}_{1-\alpha}^{(N)} = (F\otimes\mathbb{I}_{r+1})\cdot B_{1-\alpha}.
\]
Noting that $B_{1-\alpha}=B_\alpha^{-1}$ because $(1-\alpha)-\frac12=-(\alpha-\frac12)$, we have:
\[
\mathcal{U}_\alpha = (F\otimes\mathbb{I})\,B_\alpha\,B_\alpha\,(F^\dagger\otimes\mathbb{I}) = (F\otimes\mathbb{I})\,D_\alpha\,(F^\dagger\otimes\mathbb{I}),
\]
where $D_\alpha=B_\alpha^2$ has diagonal blocks $\exp(2i(\alpha-\frac12)\operatorname{Li}(x_n)\mathbf{H})$.

\textit{Sufficiency}: If $\alpha=1/2$, then $D_{1/2}=\mathbb{I}$, so $\mathcal{U}_{1/2}=(F\otimes\mathbb{I})(F^\dagger\otimes\mathbb{I})=\mathbb{I}$.

\textit{Necessity}: If $\alpha\neq 1/2$, since $\operatorname{Li}(x_n)$ is strictly monotone in $n$ and $\mathbf{H}\neq 0$ (at least one $h_j\neq 0$), the diagonal matrix $D_\alpha$ contains nontrivial phase variation. After conjugation by $F$, $\mathcal{U}_\alpha$ is no longer the identity (its off-diagonal blocks arise from non-zero interference produced by the variation of $\operatorname{Li}(x_n)$). Specifically, taking $k=0$:
\[
[\mathcal{U}_\alpha]_{0,0}=\frac{1}{N}\sum_{n=0}^{N-1}\exp(2i(\alpha-\tfrac12)\operatorname{Li}(x_n)\mathbf{H})\neq\mathbb{I}_{r+1},
\]
because the summand is non-constant in $N$.
\end{proof}

\begin{lemma}[Involution Property]\label{lem:involution}
The duality composite operator satisfies the involution relation:
\[
\mathcal{U}_\alpha^{(N)}\cdot\mathcal{U}_{1-\alpha}^{(N)}=\mathbb{I}.
\]
\end{lemma}

\begin{proof}
\[
\mathcal{U}_\alpha\cdot\mathcal{U}_{1-\alpha}=\mathcal{F}_\alpha\mathcal{F}_{1-\alpha}^{\dagger}\cdot\mathcal{F}_{1-\alpha}\mathcal{F}_\alpha^{\dagger}=\mathcal{F}_\alpha(\mathcal{F}_{1-\alpha}^{\dagger}\mathcal{F}_{1-\alpha})\mathcal{F}_\alpha^{\dagger}=\mathcal{F}_\alpha\mathbb{I}\mathcal{F}_\alpha^{\dagger}=\mathbb{I}
\]
using individual unitarity $\mathcal{F}_{1-\alpha}^{\dagger}\mathcal{F}_{1-\alpha}=\mathbb{I}$.
\end{proof}

\begin{corollary}\label{cor:unitary}
$\mathcal{U}_\alpha$ is unitary with eigenvalues on the unit circle.
\end{corollary}

\subsection{The Duality-Deviation Generator}

Since $\mathcal{U}_\alpha$ is unitary, define its logarithm (continuous branch):
\[
H_\alpha^{(N)} := -i \log \mathcal{U}_\alpha^{(N)}.
\]
From the spectral decomposition $\mathcal{U}_\alpha = F D_\alpha F^\dagger$ in Theorem~\ref{thm:selfdual}, the eigenvalues are read off directly:
\[
\theta_{n,j} = 2(\alpha-\tfrac12)\operatorname{Li}(x_n)h_j,\quad n=0,\dots,N-1,\;j=1,\dots,r+1.
\]
Their extremal behavior:
\begin{itemize}
\item Minimum ($n=0$, $\operatorname{Li}(x_0)=\operatorname{Li}(2)\approx 1.045$): $|\theta_{\min}| = 2|\alpha-\frac12|\operatorname{Li}(2)h_{\min}=O(|\alpha-\frac12|)$;
\item Maximum ($n=N-1$, $\operatorname{Li}(x_{N-1})\sim\frac{N}{\log N}$): $|\theta_{\max}| \sim 2|\alpha-\frac12|\frac{N}{\log N}h_{\max}=O(|\alpha-\frac12|\frac{N}{\log N})$.
\end{itemize}
Therefore $\|H_\alpha^{(N)}\|_{\mathrm{op}}\sim|\alpha-1/2|\cdot\frac{N}{\log N}$.

\begin{definition}[Scaled Generator]\label{def:scaled}
\[
\mathcal{H}_\alpha^{(N)} := \frac{N}{\log N} H_\alpha^{(N)}.
\]
\end{definition}

Its eigenvalues $\tilde{\theta}_{n,j} = \frac{N}{\log N}\theta_{n,j}$ satisfy:
\begin{itemize}
\item $|\tilde{\theta}_{\min}| \sim |\alpha-\frac12|\frac{N}{\log N}\operatorname{Li}(2)h_{\min}=O(\frac{N}{\log N})$;
\item $|\tilde{\theta}_{\max}| \sim |\alpha-\frac12|\frac{N^2}{(\log N)^2}h_{\max}=O(\frac{N^2}{(\log N)^2})$.
\end{itemize}

\section{Scaled Generator and Spectral Reduction to Jacobi Operators}

\subsection{Explicit Structure of Eigenvalues}

\begin{theorem}[Eigenvalue Diagonalization]\label{thm:eigendiagonal}
The duality composite operator $\mathcal{U}_\alpha^{(N)}$ is unitarily similar to the diagonal phase matrix:
\[
\mathcal{U}_\alpha^{(N)} = (F\otimes\mathbb{I}_{r+1})\,D_\alpha\,(F^\dagger\otimes\mathbb{I}_{r+1}),
\]
where $D_\alpha$ is diagonal with entries $\lambda_{n,j} = \exp(2i(\alpha-\frac12)\operatorname{Li}(x_n)h_j)$.
\end{theorem}

\begin{proof}
The factorization was already given in the proof of Theorem~\ref{thm:selfdual}. Since $F\otimes\mathbb{I}$ is unitary, this decomposition directly yields the spectrum of $\mathcal{U}_\alpha$.
\end{proof}

\begin{corollary}[Logarithmic Eigenvalues]\label{cor:logeigen}
The scaled generator $\mathcal{H}_\alpha^{(N)}$ has eigenvalues
\[
\tilde{\theta}_{n,j} = \frac{N}{\log N} \cdot 2(\alpha-\tfrac{1}{2})\operatorname{Li}(x_n)h_j.
\]
\end{corollary}

\subsection{Strict Lower Bound on Spectral Divergence}

\begin{theorem}[Unboundedness at Non-Criticality]\label{thm:unbounded}
When $\alpha \neq 1/2$, for all $n,j$:
\[
|\tilde{\theta}_{n,j}| \geq \frac{N}{\log N} \cdot 2|\alpha-\tfrac{1}{2}| \cdot \operatorname{Li}(2) \cdot h_{\min} = O\left(\frac{N}{\log N}\right),
\]
and $|\tilde{\theta}_{n,j}| \leq C \cdot |\alpha-\frac{1}{2}| \cdot \frac{N^2}{(\log N)^2}$. Hence all eigenvalues diverge uniformly as $N \to \infty$.
\end{theorem}

\begin{proof}
By strict monotonicity of $\operatorname{Li}(x)$ for $x\geq 2$, we have $\operatorname{Li}(x_n)\geq\operatorname{Li}(2)>0$ for all $n\geq 0$. Substituting into Corollary~\ref{cor:logeigen} gives the lower bound. For the upper bound, use the asymptotic formula $\operatorname{Li}(x)\sim x/\log x$ and $x_{N-1}=2e^{(N-1)\delta}\approx N$ (with $\delta=\log(N/2)/N$) to obtain:
\[
\operatorname{Li}(x_{N-1})\sim\frac{N}{\log N},\quad |\tilde{\theta}_{\max}|\sim\frac{N}{\log N}\cdot 2|\alpha-\tfrac12|\cdot\frac{N}{\log N}\cdot h_{\max}.
\]
As $N\to\infty$, the factor $\frac{N}{\log N}\to\infty$, while $|\alpha-\frac12|>0$ is fixed, so $|\tilde{\theta}_{n,j}|\to\infty$ uniformly for all $n,j$.
\end{proof}

\begin{definition}[Empirical Spectral Measure]\label{def:esm}
\[
\mu_\alpha^{(N)} := \frac{1}{N(r+1)} \sum_{j=1}^{N(r+1)} \delta_{\tilde{\theta}_j}.
\]
\end{definition}

\begin{theorem}[Weak Convergence of Spectral Measure and Phase Transition]\label{thm:spectralmeasure}
\begin{itemize}
\item When $\alpha = \frac{1}{2}$: $\mathcal{H}_{1/2}^{(N)} = 0$, so for all $N$, $\mu_{1/2}^{(N)} = \delta_0$ (Dirac measure concentrated at the origin), corresponding to compactification in the real-direction.
\item When $\alpha \neq \frac{1}{2}$: The empirical measure $\mu_\alpha^{(N)}$ satisfies:
\[
\operatorname{supp}(\mu_\alpha^{(N)}) \subseteq \left[-C|\alpha-1/2|\frac{N^2}{(\log N)^2},\; C|\alpha-1/2|\frac{N^2}{(\log N)^2}\right],
\]
and as $N \to \infty$ the support interval expands to infinity; thus $\mu_\alpha^{(N)}$ is non-compact in the limit.
\end{itemize}
\end{theorem}

\subsection{Lanczos Tridiagonalization and Jacobi Coefficients}

Applying Lanczos iteration to $\mathcal{H}_\alpha^{(N)}$ with initial vector $v_0 = \frac{1}{\sqrt{N(r+1)}}(1,\dots,1)^T$, we obtain the tridiagonal Jacobi matrix $J_\alpha^{(N)}$, where $M = N(r+1)$.

\begin{theorem}[Critical Behavior of Jacobi Coefficients]\label{thm:jacobi}
\begin{itemize}
\item When $\alpha = 1/2$: For all $n$, $\beta_n^{(N)} = 0$, $\alpha_n^{(N)} = 0$; the Jacobi matrix is zero.
\item When $\alpha \neq 1/2$: For each fixed $n$, as $N \to \infty$:
\[
\beta_n^{(N)} \sim |\alpha-1/2| \cdot \frac{N^2}{(\log N)^2} \cdot c_n \to \infty,
\]
where $c_n>0$ is a constant depending on $n$.
\end{itemize}
\end{theorem}

\begin{proof}
For $\alpha=1/2$, $\mathcal{H}_{1/2}^{(N)}=0$; the first Lanczos step gives $\alpha_0=v_0^\dagger\mathcal{H}v_0=0$, and subsequently $\beta_1=\|\mathcal{H}v_0-\alpha_0 v_0\|=0$, terminating the iteration.

For $\alpha\neq 1/2$, use the spectral representation $\mathcal{H}_\alpha^{(N)}=\sum_{j=1}^M \tilde{\theta}_j P_j$ ($P_j$ spectral projections). The initial vector $v_0$ has non-zero projection in every eigen-direction (since $v_0$ is uniform and the eigenbasis is the Fourier basis). Compute:
\[
\alpha_0 = \frac{1}{M}\sum_{j=1}^M \tilde{\theta}_j,\quad \|\mathcal{H}_\alpha v_0\|^2 = \frac{1}{M}\sum_{j=1}^M \tilde{\theta}_j^2.
\]
By Theorem~\ref{thm:unbounded}, both the maximum and minimum $\tilde{\theta}_j$ diverge, and positive/negative values are symmetrically distributed (because $\mathbf{H}$ is trace-zero, so the eigenvalues $h_j$ come in positive/negative pairs). Thus $\alpha_0$ is of the same order as $\tilde{\theta}_{\max}$, and:
\[
\beta_1^2 = \|\mathcal{H}_\alpha v_0\|^2 - \alpha_0^2 = \frac{1}{M}\sum_j\tilde{\theta}_j^2 - \left(\frac{1}{M}\sum_j\tilde{\theta}_j\right)^2.
\]
This variance is of the same order as $\tilde{\theta}_{\max}^2$, i.e., $O((\alpha-\frac12)^2\frac{N^4}{(\log N)^4})$, so $\beta_1\sim|\alpha-\frac12|\frac{N^2}{(\log N)^2}$. For higher-order $\beta_n$, by continuity of the Lanczos iteration the scaling behavior is the same as $\beta_1$, differing only in the constant $c_n$.
\end{proof}

\subsection{Carleman Condition and Essential Self-Adjointness}

\begin{theorem}[Critical Regularization of the Carleman Condition]\label{thm:carleman}
For fixed $N$, define the regularized Carleman sum
\[
S_N(\alpha) := \sum_{n=1}^{M-1} \frac{1}{\beta_n^{(N)}(\alpha)} \quad (\alpha \neq \tfrac12,\; M=N(r+1)).
\]
By Theorem~\ref{thm:jacobi}, as $\alpha \to 1/2$:
\[
S_N(\alpha) \sim \frac{(\log N)^2}{|\alpha-\tfrac12|N^2}\sum_{n=1}^{M-1}\frac{1}{c_n} = \frac{C_N}{|\alpha-\tfrac12|} \to +\infty,
\]
where $C_N>0$. Hence, at the critical parameter $\alpha=1/2$, the Carleman condition is satisfied with infinite strength. Conversely, for any fixed $\alpha \neq 1/2$:
\[
\lim_{N\to\infty} S_N(\alpha) = 0,
\]
because $\beta_n^{(N)} \sim |\alpha-\tfrac12| \frac{N^2}{(\log N)^2} c_n \to \infty$.
\end{theorem}

\textbf{Physical Interpretation}:
\begin{itemize}
\item $\alpha=1/2$ corresponds to the ``ground state'': the system sits at zero energy, no excitations, and the spectrum is extremely compactified (a single point). The Carleman condition is satisfied at a rate $1/|\alpha-1/2|$, marking the algebraic stability of the critical line.
\item $\alpha \neq 1/2$ corresponds to an ``excited state'': the level spacing $\beta_n$ grows to infinity with $N$, the Carleman sum tends to zero, and the system escapes the compact binding into a non-compact extended state. This strictly corresponds to the ``condensation'' of Riemann zeros on the critical line versus their ``divergence'' when departing from it.
\end{itemize}

\subsection{Spectral Correspondence with Riemann Zeros}

The scaled eigenvalues $\tilde{\theta}_{n,j}$ are related to the imaginary parts $\gamma_j$ of Riemann zeros (where $\rho_j = 1/2 + i\gamma_j$) through the following correspondence: under appropriate scale normalization, $\tilde{\theta}_j$ is asymptotically equal to $2\pi\gamma_j$. Specifically, the discrete frequencies $t_k = \frac{2\pi k}{\log N}$ correspond to sampling the imaginary parts of $\zeta$, while the eigenvalues $\tilde{\theta}_j$ of $\mathcal{H}_\alpha^{(N)}$ at $\alpha \to 1/2$ are related to $t_k$ via the scale factor $\frac{N}{\log N}$; combining this with $\operatorname{Li}(x_n) \sim \frac{N}{\log N}$ yields the above asymptotic correspondence. Therefore, the contraction behavior of the empirical measure $\mu_\alpha^{(N)}$ as $\alpha \to 1/2$ corresponds to the ``condensation'' phenomenon of Riemann zeros on the critical line $\Re(s)=1/2$.

\section{DFT/FFT and the Riemann Functional Equation}

\subsection{Poisson Summation and Duality}

The standard DFT is a discretization of the Poisson summation formula:
\[
\sum_{n\in\mathbb{Z}} f(n) = \sum_{k\in\mathbb{Z}} \hat{f}(k), \quad \hat{f}(\xi) = \int_{-\infty}^\infty f(x) e^{-2\pi i x\xi}\,dx.
\]
For suitable test functions $f\in\mathcal{S}(\mathbb{R}^+)$ (rapidly decreasing), the Mellin transform gives
\[
\int_0^\infty x^{s-1} e^{-2\pi i k x}\, dx = (2\pi i k)^{-s} \Gamma(s).
\]
After discretization, the DFT matrix element $\frac{1}{\sqrt{N}} e^{2\pi i k n/N}$ corresponds to sampling the multiplicative character.

\subsection{Operator Realization of the Riemann Functional Equation}

\begin{theorem}[Functional Equation Operator Realization]\label{thm:operatorfe}
In the weak topology on the space of rapidly decreasing functions $\mathcal{S}(\mathbb{R})$ (equivalently, distributional convergence in $\mathcal{S}'(\mathbb{R})$),
\[
\mathcal{U}_\alpha^{(N)} \xrightarrow{w} \mathcal{M}_{\chi(\alpha)/|\chi(\alpha)|}, \quad (N\to\infty)
\]
where $\mathcal{M}_{g}$ denotes the multiplication operator $(\mathcal{M}_g \hat{f})(t) = g(t)\hat{f}(t)$, and $\chi(s)=2^s\pi^{s-1}\sin(\frac{\pi s}{2})\Gamma(1-s)$ is the factor appearing in the functional equation of Riemann $\zeta$.
\end{theorem}

\begin{proof}[Proof Framework]
Use the Mellin--Plancherel method. The NCDFT kernel $K_{\alpha}^{(N)}(x_{n},t_{k})= x_{n}^{\alpha-1/2}e^{-i t_{k}\log x_{n}}$ approximates the continuous kernel $x^{\alpha-1/2+it}$ as $N \to\infty$. For any test function $\varphi\in\mathcal{S}(\mathbb{R})$, consider the bilinear form:
\[
\langle \mathcal{U}_\alpha^{(N)}\hat{\varphi},\hat{\varphi}\rangle = \frac{1}{N}\sum_{k,k'}\overline{\hat{\varphi}(t_k)}\hat{\varphi}(t_{k'})\sum_{n=0}^{N-1} e^{2\pi i(k-k')n/N} \exp\left(2i(\alpha-\tfrac12)\operatorname{Li}(x_n)\mathbf{H}\right).
\]
Writing the inner sum as a Riemann sum and using Poisson summation with Plancherel's theorem for Mellin transforms, the limit becomes:
\[
\int \overline{\hat{\varphi}(t)}\,\frac{\chi(\alpha+it)}{|\chi(\alpha+it)|}\,\hat{\varphi}(t)\,dt.
\]
The phase normalization $\chi/|\chi|$ ensures boundedness of the limit operator (modulus 1), consistent with the unitarity of $\mathcal{U}_\alpha$. When $\alpha=1/2$, $\chi(1/2+it)/|\chi(1/2+it)|=1$, so $\mathcal{U}_{1/2}=\mathbb{I}$ consistent with Theorem~\ref{thm:selfdual}.
\end{proof}

\begin{center}
\textbf{Riemann Functional Equation Correspondence Table}
\medskip
\begin{tabular}{|l|l|l|}
\hline
Riemann $\zeta$ function & NCDFT operator & Limit behavior ($N\to\infty$) \\
\hline
Complex parameter $s = \sigma + it$ & Real parameter $\alpha = \sigma$ (real part) + frequency index $k\leftrightarrow t$ (imaginary part) & At $\alpha=1/2$, $t_k \to \gamma_j$ \\
Symmetry axis $\Re(s) = 1/2$ & Critical value $\alpha = 1/2$ ($\mathcal{H}_{1/2}^{(N)}=0$) & Real part locked, imaginary part expands as discrete set \\
Functional equation $s \leftrightarrow 1-s$ & Duality pair $\mathcal{F}_\alpha \leftrightarrow \mathcal{F}_{1-\alpha}$ & At $\alpha=1/2$, $\mathcal{U}=\mathbb{I}$ (self-dual) \\
Self-duality on critical line $\zeta(1/2+it) \leftrightarrow \overline{\zeta(1/2+it)}$ & Self-duality $\mathcal{U}_{1/2} = \mathbb{I}$ & Phase normalized to 1 \\
\hline
\end{tabular}
\end{center}

\subsection{From Input Function to the Logarithmic Derivative of $\zeta$}

Let $f_n = \frac{\Lambda(n)}{\sqrt{n}} \phi\left(\frac{\log n}{\log N}\right)$, where $\phi$ is as defined in Section~1.1 (support $[0,1]$). Consider its DFT:
\[
\hat{f}_k = \frac{1}{\sqrt{N}}\sum_{n=1}^{N-1} f_n e^{-2\pi i k n/N}.
\]

\begin{theorem}[Smooth Truncation Correspondence]\label{thm:smoothtrunc}
For $t_k = \frac{2\pi k}{\log N}$ ($|k|\le N$), we have
\[
\frac{1}{\sqrt{N}}\hat{f}_k = -\frac{\zeta'(1/2+it_k)}{\zeta(1/2+it_k)}\cdot\frac{\log N}{\sqrt{N}} + \mathcal{E}_N(t_k),
\]
where the error $\mathcal{E}_N(t)$ satisfies the $L^2$ average estimate
\[
\frac{1}{N}\sum_{k=0}^{N-1} |\mathcal{E}_N(t_k)|^2 = O(N^{-1+2\epsilon}) \quad (\forall \epsilon>0).
\]
\end{theorem}

\begin{proof}[Proof Sketch]

\textit{Step 1: Mellin-transform representation.} Let $\Phi_N(x)=\phi(\frac{\log x}{\log N})$, whose Mellin transform is $\Phi_N^*(u)=\log N \int_0^1 \phi(y) N^{uy}\,dy$. Then
\[
S_N(t):=\sum_{n=1}^\infty \frac{\Lambda(n)}{n^{1/2+it}}\Phi_N(n) = \frac{1}{2\pi i}\int_{(c)} \Phi_N^*(u)\left(-\frac{\zeta'(1/2+it+u)}{\zeta(1/2+it+u)}\right)du,
\]
with integration path $\Re(u)=c>1/2$.

\textit{Step 2: Contour shift and residues.} Move the contour to $\Re(u)=-\delta$ ($0<\delta<1/2$), crossing the pole at $u=0$ (simple pole of $-\zeta'/\zeta$, residue $-\frac{\zeta'}{\zeta}(1/2+it)$) and the zero poles at $u=\rho-1/2-it$. The main term comes from $u=0$:
\[
-\frac{\zeta'}{\zeta}(1/2+it)\cdot\Phi_N^*(0) = -\frac{\zeta'}{\zeta}(1/2+it)\log N,
\]
since $\Phi_N^*(0)=\log N\int_0^1\phi(y)\,dy=\log N$.

\textit{Step 3: Zero residue sum.} Using the exponential decay $\hat{\phi}(-iy)\sim e^{-2\pi y}$ (since $\operatorname{supp}\phi\subset[0,1]$), combined with Ingham's zero-density estimate $N(\sigma,T)\ll T^{4(1-\sigma)+\epsilon}$, off-critical zeros ($\Re\rho\neq1/2$) contribute $O(N^{-c})$; critical-line zeros contribute via local density $O(\log N)$, absorbed by averaging.

\textit{Step 4: Horizontal integrals.} Using growth estimates for $-\zeta'/\zeta$, horizontal contributions are $O(N^{-\delta+\epsilon})$.

\textit{Step 5: Discretization error.} Replacing $e^{-2\pi i k n/N}$ by $e^{-i t_k\log n}$ differs by $O(N^{-1})$ on average via Cauchy--Schwarz and the Prime Number Theorem. Combining yields the theorem.
\end{proof}

\subsection{Sampling Completeness and Zero Approximation}

\begin{lemma}[Grid Density]\label{lem:density}
Let $\rho_j = \frac{1}{2} + i\gamma_j$ be the non-trivial zeros of the Riemann $\zeta$ function (ordered by imaginary part $0 < \gamma_1 \leq \gamma_2 \leq \cdots$). The discrete frequency grid is defined by:
\[
t_k = \frac{2\pi k}{\log N}, \quad k = 0, 1, \dots, N-1.
\]
Then for each fixed zero index $j$, as $N \to \infty$:
\[
\lim_{N\to\infty} \min_{0 \leq k < N} |t_k - \gamma_j| = 0.
\]
The approximation error order is:
\[
\min_k |t_k - \gamma_j| \leq \frac{\pi}{\log N} = O\left(\frac{1}{\log N}\right).
\]
\end{lemma}

\begin{proof}
The sampling step is $\Delta t = t_{k+1} - t_k = \frac{2\pi}{\log N}$, which tends to $0$ monotonically as $N\to\infty$. By the Riemann--von Mangoldt formula, the zeros $\gamma_j$ are finite and isolated, with minimum spacing $\delta_j = \min_{j'\neq j}|\gamma_{j'}-\gamma_j| > 0$. Take $N$ large enough that $\frac{2\pi}{\log N} < \frac{\delta_j}{2}$; then there must exist an integer $k$ satisfying $|t_k - \gamma_j| \leq \frac{\Delta t}{2} = \frac{\pi}{\log N}$.
\end{proof}

\begin{proposition}[FFT Pole-Zero Correspondence]\label{prop:polezero}
Near a zero $\gamma_j$, the FFT amplitude $|D_N(t_k)|$ exhibits Lorentzian resonance behavior:
\[
|D_N(t_k)|^2 \sim \frac{1}{(t_k - \gamma_j)^2 + \epsilon_N^2}, \quad \text{when } |t_k - \gamma_j| \ll 1,
\]
where the effective resolution width $\epsilon_N \sim (\log N)^{-1}$ is determined jointly by the sampling density and truncation error. Therefore, local maxima converge to the true zero locations:
\[
\arg\max_{t_k} |D_N(t_k)| \xrightarrow{N\to\infty} \{\gamma_j\}_{j=1}^\infty.
\]
\end{proposition}

\begin{proof}[Proof Sketch]
By Theorem~\ref{thm:smoothtrunc}, $D_N(t_k)$ approximates $-\frac{\zeta'}{\zeta}(1/2+it_k)$ in the $L^2$ sense. The latter has a simple pole at $\gamma_j$:
\[
-\frac{\zeta'(1/2+it)}{\zeta(1/2+it)} \sim \frac{1}{i(t-\gamma_j)}\quad (t\to\gamma_j),
\]
so the amplitude diverges near $t\approx\gamma_j$. Combined with Lemma~\ref{lem:density} (grid refinement), there must exist a subsequence $t_{k(N)}\to\gamma_j$ such that $|D_N(t_{k(N)})|\to\infty$ (relative to neighboring points), forming a detectable peak.
\end{proof}

\begin{corollary}[Phase Jump Criterion]\label{cor:phasejump}
At discrete points near a zero, the FFT phase undergoes a $\pi$ jump:
\[
\arg(D_N(t_{k_j+1})) - \arg(D_N(t_{k_j-1})) \approx \pi \cdot \mathrm{sgn}(t_{k_j} - \gamma_j) + O\left(\frac{1}{\log N}\right),
\]
where $k_j = \arg\min_k |t_k - \gamma_j|$. This criterion is unaffected by amplitude normalization and is a more robust numerical detection standard.

\textbf{Scope of validity}: Theorem~\ref{thm:smoothtrunc}'s $L^2$ error estimate is valid for $t \in [0, T(N)]$ with $T(N) = o(\sqrt{N})$. Therefore, to detect the first $J$ zeros, one needs $\gamma_J < \sqrt{N}$, i.e., $J \log J \lesssim \sqrt{N}$. This indicates that $N$ must grow quadratically with the number of zeros to be detected, consistent with the constructive protocol of Section~7.3.
\end{corollary}

\section{Operator Realization of the Weil Formula}

\subsection{Complex Spectral Lift Operator and Double Measure Structure}

From Section~3, the scaled generator $\mathcal{H}_\alpha^{(N)}$ has eigenvalues $\tilde{\theta}_j$ asymptotically corresponding to the imaginary parts $\gamma_j$ of Riemann zeros as $\alpha\to1/2,N\to\infty$. Define the spectral lift operator:
\[
\mathcal{L}_\alpha^{(N)} := \frac{1}{2}\mathbb{I} + i\frac{\mathcal{H}_\alpha^{(N)}}{2\pi}.
\]

\textbf{Note}: When strictly at $\alpha=1/2$, $\mathcal{H}_{1/2}^{(N)}=0$ causes the ``internal'' eigenvalues of $\mathcal{L}_{1/2}^{(N)}$ to be the real number $1/2$. The full spectral lift must incorporate the DFT frequency indices $t_k$, understood as:
\[
\mathcal{L}_{1/2}^{\mathrm{full},(N)} = \frac{1}{2}\mathbb{I}_{N(r+1)} + i\, \mathrm{diag}(t_0,\dots,t_{N-1}) \otimes \mathbb{I}_{r+1},
\]
whose eigenvalues are $\frac{1}{2} + i t_k$, approximating the zeros $\frac{1}{2}+i\gamma_j$ as $N\to\infty$.

\textbf{Double measure structure}:

\begin{enumerate}
\item \textbf{Internal spectral measure $\mu_\alpha^{(N)}$ (real-part control)}:
\[
\mu_\alpha^{(N)} := \frac{1}{N(r+1)}\sum_{j=1}^{N(r+1)} \delta_{\tilde{\theta}_j}.
\]
\begin{itemize}
\item $\alpha=1/2$: $\mathcal{H}_{1/2}=0$, so $\mu_{1/2}^{(N)} = \delta_0$, indicating that the real part is strictly locked at $1/2$.
\item $\alpha\neq 1/2$: support diverges, Carleman condition fails.
\end{itemize}

\item \textbf{Frequency spectral measure $\nu_{1/2}^{(N)}$ (imaginary-part expansion)}:
\[
\nu_{1/2}^{(N)} := \frac{1}{\log N} \sum_{k=0}^{N-1} |D_N(t_k)|^2 \,\delta_{\frac{1}{2} + i t_k}.
\]
This measure is always supported on the critical line $\Re(s)=1/2$.
\end{enumerate}

\textbf{Key supplement}: The support of this measure is $\{\frac{1}{2} + i t_k \mid k=0,1,\dots,N-1\}$. By Lemma~\ref{lem:density} (sampling completeness), as $N\to\infty$ these discrete points become uniformly dense on the critical line $\Re(s)=1/2$, and for every zero $\rho_j$ there exists a subsequence $t_{k(j,N)} \to \gamma_j$. This provides the geometric foundation for the weak convergence of the measure.

The duality $\mathcal{L}_{1-\alpha} = \frac{1}{2} - i\mathcal{H}_\alpha \cdot \frac{1}{2\pi} = \overline{\mathcal{L}_\alpha}$ embodies the complex-conjugate symmetry of the functional equation $s\leftrightarrow 1-s$.

\subsection{Arithmetic-Side Representation and Joint Limit Theorem}

To avoid direct trace divergence, use the empirical spectral measure (normalized trace):

For $\mu_\alpha^{(N)}$:
\[
\mu_\alpha^{(N)}[F] := \frac{1}{N(r+1)}\operatorname{Tr} F(\mathcal{L}_\alpha^{(N)}) = \frac{1}{N(r+1)}\sum_{j=1}^{N(r+1)} F\left(\frac{1}{2}+i\frac{\tilde{\theta}_j}{2\pi}\right).
\]

For $\nu_{1/2}^{(N)}$:
\[
\nu_{1/2}^{(N)}[F] := \frac{1}{\log N}\sum_{k=0}^{N-1} F\left(\frac{1}{2}+it_k\right) \cdot |D_N(t_k)|^2.
\]

On the other hand, for smooth compactly supported functions $F$, consider its inverse Mellin transform $\hat{F}(x)=\frac{1}{2\pi i}\int_{(c)} F(s) x^{-s}\,ds$. Using Parseval's identity and Theorem~\ref{thm:smoothtrunc}, we obtain:
\[
\nu_{1/2}^{(N)}[\widehat{F}] = \frac{1}{\log N}\sum_{n=1}^\infty \frac{\Lambda(n)}{\sqrt{n}}\phi\left(\frac{\log n}{\log N}\right)\widehat{F}(\log n) + o(1),
\]
where $\widehat{F}(t) = \int_{-\infty}^\infty F(1/2+i\tau) e^{it\tau}\, d\tau$ is the Fourier transform of $F$.

\begin{theorem}[Convergence at Fixed $\alpha=1/2$]\label{thm:weilconv}
Let $\nu_{1/2}^{(N)}$ be the weighted empirical measure constructed from the discrete frequencies $t_k = \frac{2\pi k}{\log N}$. Then for any smooth compactly supported test function $F$:
\[
\lim_{N\to\infty} \nu_{1/2}^{(N)}[\widehat{F}] = \frac{1}{2\pi} \sum_{\rho} F(\rho),
\]
where $\rho$ ranges over all non-trivial zeros of the Riemann $\zeta$ function.
\end{theorem}

\begin{proof}
By Theorem~\ref{thm:smoothtrunc}, $|D_N(t_k)|^2$ approximates $|-\zeta'/\zeta(1/2+it_k)|^2$ in the $L^2$ average sense. Combining the Poisson summation formula with the sampling completeness of Lemma~\ref{lem:density}, apply the Weil explicit formula:
\[
\sum_{\rho} F(\rho) = \int_1^\infty \left(\psi(x)-x\right)\widehat{F}(\log x)\frac{dx}{x} + \text{trivial contributions},
\]
where $\psi(x)=\sum_{n\leq x}\Lambda(n)$. Transform the integral via Abel summation into a sum over $\Lambda(n)$, matching the arithmetic representation of $\nu_{1/2}^{(N)}$. Standard Tauberian arguments (using the smoothness of $\phi$ and rapid decay of $\widehat{F}$) control the truncation error to $o(1)$, finally yielding weak convergence of the measure to the zero distribution, up to the factor $1/(2\pi)$ from Fourier-transform normalization.
\end{proof}

\subsection{Constructive Repulsion of Off-Axis Zeros}\label{sec:repulsion}

\begin{proposition}[Constructive Repulsion of Off-Axis Zeros]\label{prop:repulsion}
Suppose there exists a non-trivial zero $\rho_0 = \sigma_0 + i\gamma_0$ with $\sigma_0 \neq 1/2$. Without loss of generality assume $\sigma_0 > 1/2$ (by symmetry of the functional equation). Consider the frequency measure $\nu_{1/2}^{(N)}$ generated by NCDFT at $\alpha=1/2$. By the Mellin--Plancherel representation of Theorem~\ref{thm:smoothtrunc}, the asymptotic expansion of $D_N(t)$ at $t=\gamma_0$ contains an explicit residue contribution from the off-axis zero:
\[
D_N(\gamma_0) = -\Phi_N^*(\sigma_0-\tfrac12) \cdot \mathop{\mathrm{Res}}_{s=\rho_0}\frac{\zeta'(s)}{\zeta(s)} + O((\log N)^2),
\]
where $\Phi_N^*(u) = \log N \int_0^1 \phi(y) N^{uy}\, dy$ is the Mellin transform of the truncation kernel. By the Laplace method (saddle-point method), for $\sigma_0 > 1/2$:
\[
\Phi_N^*(\sigma_0-\tfrac12) = \log N \int_0^1 \phi(y) e^{(\sigma_0-1/2)y\log N}\, dy.
\]
Since $\phi(1)>0$ and $\phi$ is non-zero near $y=1$, this integral is dominated by the endpoint $y=1$:
\[
\Phi_N^*(\sigma_0-\tfrac12) \sim c_\phi \frac{N^{\sigma_0-1/2}}{\sigma_0-1/2} \quad (N\to\infty),
\]
where $c_\phi = \phi(1) > 0$ (guaranteed by Section~1.1). Therefore:
\[
|D_N(\gamma_0)| \sim C_{\phi,\rho_0} \cdot N^{\sigma_0-1/2}.
\]
Meanwhile, the peak amplitude produced by critical-line zeros $\rho_j = 1/2+i\gamma_j$ on the discrete grid $t_k$ is only $O((\log N)^2)$ (Proposition~\ref{prop:polezero}). Thus the amplitude corresponding to the off-axis zero exponentially overwhelms the critical-line zeros by a power of $N$.

This means the frequency measure $\nu_{1/2}^{(N)}$ would be forced to carry an off-axis component of weight
\[
|D_N(t_{k(N)})|^2 \sim C_{\phi,\rho_0}^2 \cdot N^{2(\sigma_0-1/2)}
\]
at grid points $t_{k(N)}$ in an $O(1/\log N)$ neighborhood of $\gamma_0$. However, by the constructive definition of $\nu_{1/2}^{(N)}$ (Section~5.1), its support strictly satisfies:
\[
\mathrm{supp}(\nu_{1/2}^{(N)}) = \left\{\frac{1}{2}+it_k \;\middle|\; k=0,\dots,N-1\right\} \subset \{s\in\mathbb{C}:\Re(s)=1/2\}.
\]
This support contains no off-axis points. Hence the existence of an off-axis zero leads to a constructive contradiction: the explicit formula requires the measure to carry exponentially growing off-axis mass near the critical line, while the NCDFT construction at $\alpha=1/2$ strictly forbids any off-axis support. The assumption is therefore false, and all non-trivial zeros must satisfy $\Re(\rho)=1/2$.
\end{proposition}

\subsection{Limit Order and Critical-Line Compactification}

\begin{theorem}[Limit Order and Zero Distribution]\label{thm:limitorder}\ 

\begin{enumerate}
\item[(A)] \textbf{Correct order (compactifying limit)}:
Fix $\alpha=1/2$, consider the measure $\nu_{1/2}^{(N)}$ as defined above; then as $N\to\infty$:
\[
\lim_{N\to\infty} \nu_{1/2}^{(N)} = \frac{1}{2\pi}\sum_{\rho} \delta_{\rho}.
\]

\item[(B)] \textbf{Wrong order (non-compact divergence)}:
If we first fix $\alpha \neq 1/2$ and let $N\to\infty$, the scaled generator satisfies $\|\mathcal{H}_{\alpha}^{(N)}\|_{\mathrm{op}} \sim |\alpha-1/2|\frac{N^2}{(\log N)^2} \to \infty$, and the spectral-measure support expands to infinity (non-compact). If we then let $\alpha\to 1/2$, the limit does not exist (or degenerates into a meaningless divergence).

\item[(C)] \textbf{Non-commutativity of limits}:

For the internal measure $\mu_\alpha^{(N)}$:
\begin{itemize}
\item First $N\to\infty$ (fixed $\alpha\neq 1/2$): $\|\mathcal{H}_\alpha^{(N)}\|\to\infty$, $\mu_\alpha^{(N)}$ has no compact support, weak limit does not exist.
\item First $\alpha\to 1/2$ (fixed $N$): yields $\delta_0$, then $N\to\infty$ remains $\delta_0$ (degenerate point measure).
\end{itemize}

For the frequency measure $\nu_{1/2}^{(N)}$ (defined only at $\alpha=1/2$):
\begin{itemize}
\item Correct order: fix $\alpha=1/2$ (lock real part), then $N\to\infty$ (expand imaginary part), yielding $\frac{1}{2\pi}\sum_\rho \delta_\rho$.
\item Reverse path: if first $N\to\infty$ then $\alpha\to 1/2$, since $\nu$ is undefined (or corresponds to non-compact operators) when $\alpha\neq 1/2$, the path is illegitimate.
\end{itemize}

\item[(D)] \textbf{Grid-refinement limit (supplement)}:
Fix $\alpha=1/2$; as $N\to\infty$, the sampling step $\Delta t = \frac{2\pi}{\log N} \to 0$, and the grid $t_k$ becomes dense on any compact interval $[0,T]$. This limit is independent of the $\alpha$ limit, but provides the geometric foundation for the weak convergence of $\nu_{1/2}^{(N)}$ (see Section~4.4).
\end{enumerate}
\end{theorem}

\subsection{Limit Order and Functorial Fidelity}

\begin{definition}[Functorial Fidelity Measure]\label{def:fidelity}
\[
d_{\mathrm{spec}}(\alpha) := \limsup_{N\to\infty} \left| \nu_\alpha^{(N)}[\widehat{F}] - \frac{1}{2\pi}\sum_{\rho}F(\rho) \right|
\]
(taking supremum over an appropriate function class $F$), where $\nu_\alpha^{(N)}$ is the frequency-weighted measure corresponding to $\alpha$ (when $\alpha=1/2$ this is $\nu_{1/2}^{(N)}$, while for $\alpha\neq1/2$ it is defined as the corresponding non-compact measure).
\end{definition}

By Theorem~\ref{thm:limitorder}:
\begin{itemize}
\item $\alpha=1/2$ yields $d_{\mathrm{spec}}(1/2)=0$.
\item $\alpha\neq1/2$ yields $d_{\mathrm{spec}}(\alpha) > 0$ (in fact, diverges).
\end{itemize}

\section{Error Bounds and Convergence Analysis}

\subsection{Error Inequalities}

Summarizing the above results into the following inequalities:

\begin{enumerate}
\item \textbf{Carleman condition}:
\[
S_N(\alpha)=\sum_{n=1}^{M-1}\frac{1}{\beta_n^{(N)}(\alpha)} \begin{cases} \sim C_N/|\alpha-1/2| \to \infty & \alpha\to 1/2 \\ \to 0 & \alpha\neq 1/2,\; N\to\infty \end{cases}
\]
where $M=N(r+1)$.

\item \textbf{Phase stability (Lipschitz)}:
\[
\| e^{i\alpha\operatorname{Li}(x)\mathbf{H}} - e^{i\alpha'\operatorname{Li}(x)\mathbf{H}}\| \le |\alpha-\alpha'|\cdot\operatorname{Li}(x)\cdot\|\mathbf{H}\|
\]

\item \textbf{Explicit-formula remainder estimate}:
\[
\left| \frac{1}{\log N}\sum_{k} |D_N(t_k)|^2 \widehat{F}(t_k) - \frac{1}{2\pi}\sum_{\rho} F(\rho) \right| \le \epsilon(N) \sup|F|,
\]
where $\epsilon(N)=O(N^{-1/2})$ (analytic error) plus $O(1/\log N)$ (sampling error), with the sampling error actually dominating.

\item \textbf{Topological acceptance rate}:
\[
\mathcal{A}(\alpha) \ge \mathcal{A}_{\max}\exp\left(-\frac{(\alpha-1/2)^2}{2\sigma^2}\right)-\delta_N
\]

\item \textbf{Bishop convergence (Cauchy sequence)}:
\[
\|\mathcal{F}_{1/2}^{(N)}-\mathcal{F}_{1/2}^{(M)}\|_{\mathrm{norm}} < 2^{-k}\quad \text{for } N,M>N(k)
\]
where $\|A\|_{\mathrm{norm}} = \|A\|_F / \sqrt{\min(N,M)}$, $N(k) = \lceil C \cdot 2^{2k} \rceil$.

\item \textbf{RH error bound}:
\[
\left| \sum_\rho \frac{x^\rho}{\rho} - \operatorname{Tr}\left(\frac{x^{\mathcal{L}_{1/2}^{(N)}}}{\mathcal{L}_{1/2}^{(N)}}\right) \right| < \epsilon(N) x^{1/2}\ln x,\quad \epsilon(N)=O(N^{-1/2})
\]
\end{enumerate}

\textbf{Error source distinction}:
\begin{itemize}
\item Analytic error $\epsilon_{\text{analytic}}(N) = O(N^{-1/2})$: from truncation $\sum_{n>N} \Lambda(n)/\sqrt{n}$, Mellin transform, and approximation by smooth truncation $\phi$.
\item Sampling error $\delta_{\text{geom}}(N) = O(1/\log N)$: from geometric deviation between the discrete grid $t_k$ and continuous zeros $\gamma_j$ (see Lemma~\ref{lem:density}).
\end{itemize}
Since $\sqrt{N} \gg \log N$, for sufficiently large $N$ the sampling error dominates zero-location accuracy, while the analytic error is relatively minor.

\subsection{Explicit Derivation of $\epsilon(N)$}

From the truncation-error analysis of Section~4.2 and the numerical stability of FFT:
\[
\epsilon(N) \leq \frac{C_1}{\sqrt{N}} + \frac{C_2}{N} + C_3 e^{-c\sqrt{\log N}}
\]
where the first term comes from truncation $\sum_{n>N} \Lambda(n)/\sqrt{n}$, the second from discretization error $\log n \approx n\delta$, and the third from the Prime Number Theorem remainder ($\psi(x)=x+O(xe^{-c\sqrt{\log x}})$). Thus $\epsilon(N) = O(N^{-1/2})$.

\subsection{Constructive Convergence Protocol}

Given precision $\epsilon > 0$, the computational protocol is:

\begin{enumerate}
\item Choose $N = O(\epsilon^{-2})$ (ensuring analytic error $<\epsilon/2$) and $N > \exp(\pi/\epsilon)$ (ensuring sampling error $<\epsilon/2$, see Section~4.4).
\item Construct $\mathcal{F}_{1/2}^{(N)}$ (standard DFT matrix, with logarithmic-scale sampling $x_n = 2e^{n\delta}$).
\item Execute FFT to compute $D_N(t_k)$.
\item Locate $t_k \approx \gamma_j$ by detecting local maxima of $|D_N(t_k)|$ (or $\pi$ phase jumps of $\arg(D_N(t_k))$). By Lemma~\ref{lem:density}, when $N > \exp(\pi/\delta)$ one can guarantee positioning error $|t_k - \gamma_j| < \delta$ for all $\gamma_j < T$.
\item Error is controlled by inequality (6): $|\rho_j - \rho_j^{\mathrm{true}}| < \epsilon$.
\end{enumerate}

\section{Numerical Verification}

\subsection{Discrete Dirichlet Polynomials}

Define the finite truncation of the logarithmic derivative (smooth truncation):
\[
D_N(t)=\sum_{n=1}^\infty \frac{\Lambda(n)}{\sqrt{n}}\phi\left(\frac{\log n}{\log N}\right) e^{-it\log n}.
\]

\subsubsection{Asymptotic Density of the Sampling Grid}

By Section~4.4, the frequency sampling points $t_k = \frac{2\pi k}{\log N}$ have spacing $\Delta t = \frac{2\pi}{\log N}\to 0$ as $N\to\infty$. According to the Riemann--von Mangoldt zero-density estimate $N(T) \sim \frac{T}{2\pi}\log T$, for the first $J$ zeros $\gamma_1,\dots,\gamma_J$, when:
\[
N > \exp\left(\frac{2\pi}{\min_{j\leq J} |\gamma_{j+1}-\gamma_j|}\right),
\]
the grid can resolve all first $J$ zeros (no omission, no aliasing).

\subsection{FFT Phase Correspondence}

At discrete points $t_k=\frac{2\pi k}{\log N}$:
\[
D_N(t_k) \approx \sqrt{N}\cdot \mathrm{FFT}[\Lambda(n)/\sqrt{n}\,\phi(\cdot)](k).
\]
Phase-error analysis:
\[
\arg(D_N(t)) - \arg\left(-\frac{\zeta'(1/2+it)}{\zeta(1/2+it)}\right) = O\left(\frac{t^2}{\log N}\right) + O\left(\frac{1}{\sqrt{N}}\right)
\]
For $t \in [0, T]$, choosing $N \sim T^2$ makes the error arbitrarily small.

\subsection{Spectral Conditions for Zero Detection}

At a zero $\rho = 1/2 + i\gamma$, the function $-\zeta'/\zeta$ has a simple pole, and the phase undergoes a $\pi$ jump.

In the FFT spectrum this manifests as amplitude minima and phase discontinuities:
\[
|\mathrm{FFT}(k)| \approx \left| \frac{1}{\rho - (1/2+it_k)} \right| \quad \text{at } t_k \approx \gamma,
\]
forming Lorentzian-type peaks (or valleys, depending on normalization).

\textbf{Theoretical guarantee of the phase-jump criterion}:
By Corollary~\ref{cor:phasejump}, when $t_k$ passes through $\gamma_j$ from left to right, $\arg(D_N(t))$ undergoes a $+\pi$ jump (conversely $-\pi$). Detection of this jump is unaffected by amplitude normalization $|D_N(t_k)|^2/\|D_N\|_2^2$ and is a more robust numerical criterion.

\textbf{Applicability}: Theorem~\ref{thm:smoothtrunc}'s $L^2$ error estimate is valid on $[0,T(N)]$ requiring $T(N) = O(\sqrt{N})$. Thus, to detect the first $J$ zeros, we need $\gamma_J < \sqrt{N}$, i.e., $J \log J \lesssim \sqrt{N}$. This indicates that $N$ must grow quadratically with the number of zeros to be detected, consistent with the constructive protocol in Section~7.3.

\section{Conclusion: Mathematical Completeness Statement}

Based on the foregoing NCDFT categorical construction, the spectral phase transition of the scaled generator, and the dual realization of DFT/FFT with the Riemann functional equation, this section presents the functorial-rigidity version of the RH equivalence theorem. This version hardens analytic duality into purely algebraic recursive normal-form uniqueness, and excludes off-axis zeros through a double defense line of constructive contradiction and limit argument, yielding the cleanest equivalence chain.

\subsection{Recursive Block-Triangularization Rigidity}

\begin{lemma}[Recursive Block-Triangularization Rigidity]\label{lem:rigid}
Let $N=2^m$. Define the phase weight:
\[
W_n^{(N)} := \exp\!\big(i(\alpha-\tfrac12)\operatorname{Li}(x_n^{(N)})\mathbf{H}\big).
\]
The NCDFT family $\{\mathcal{F}_\alpha^{(N)}\}_{N\geq 1}$ admits a recursive block-triangularization compatible with the dyadic arithmetic decomposition $\mathbb{Z}/N\mathbb{Z} \cong \mathbb{Z}/(N/2)\mathbb{Z} \times \{0,1\}$ (i.e., a Cooley--Tukey-type butterfly recursion, reducing $\mathcal{F}_\alpha^{(N)}$ stepwise to size-$N/2$ subproblems) \textbf{if and only if} $\alpha=1/2$.
\end{lemma}

\begin{proof}
For $\alpha=1/2$, we have $W_n\equiv\mathbb{I}_{r+1}$, and $\mathcal{F}_{1/2}^{(N)}$ reduces to the standard block DFT. The standard FFT butterfly operation directly yields the block-diagonal normal form, and the recursive chain is complete.

For $\alpha\neq 1/2$, examine the logarithmic-scale sampling $x_n^{(N)}=2e^{n\delta_N}$ with $\delta_N=\frac{\log(N/2)}{N}$. The dyadic recursion requires that even-column weights agree with half-size problem weights under block permutation:
\[
W_{2m}^{(N)} = W_m^{(N/2)}.
\]
But
\[
x_{2m}^{(N)} = 2\exp\!\left(\frac{2m\log(N/2)}{N}\right),\qquad 
x_m^{(N/2)} = 2\exp\!\left(\frac{2m\log(N/4)}{N}\right).
\]
Since $\log(N/2)\neq\log(N/4)$, we have $x_{2m}^{(N)}\neq x_m^{(N/2)}$, and therefore:
\[
W_{2m}^{(N)} = \exp\!\big(i(\alpha-\tfrac12)\operatorname{Li}(x_{2m}^{(N)})\mathbf{H}\big) \neq 
\exp\!\big(i(\alpha-\tfrac12)\operatorname{Li}(x_m^{(N/2)})\mathbf{H}\big) = W_m^{(N/2)}.
\]
Thus the ``even part'' and ``odd part'' after the butterfly operation cannot be recursively reduced to an NCDFT of size $N/2$; the recursive block-triangularization chain breaks at the first step.
\end{proof}

\begin{remark}
The ``block-triangularization'' here is not a static block LU decomposition at a single $N$, but requires the normal form to evolve functorially with the dyadic arithmetic structure of $N$. It is precisely this $N$-consistency (global section) requirement that uniquely pins down $\alpha=1/2$.
\end{remark}

\subsection{NCDFT--RH Equivalence (Functorial-Rigidity Version)}

\begin{theorem}[NCDFT--RH Equivalence, Functorial-Rigidity Version]\label{thm:main}
The following statements are equivalent:
\begin{enumerate}
\item \textbf{Riemann Hypothesis}: All non-trivial zeros satisfy $\Re(\rho)=1/2$.
\item \textbf{Recursive normal form}: The NCDFT block-weighted Vandermonde matrix family $\{\mathcal{F}_\alpha^{(N)}\}_{N\geq 1}$ admits a recursive block-triangularization normal form compatible with dyadic arithmetic structure.
\item \textbf{Parameter uniqueness}: This recursive normal form uniquely corresponds to $\alpha=1/2$, and at this point the scaled generator $\mathcal{H}_{1/2}^{(N)}\equiv 0$.
\item \textbf{Spectral corollary}: The zero-distribution measure, as the spectral image of this unique section, has its support strictly collinear with $\Re(s)=1/2$.
\end{enumerate}
\end{theorem}

\begin{proof}
\begin{itemize}
\item \textbf{(2) $\Leftrightarrow$ (3)}: By Lemma~\ref{lem:rigid} (recursive block-triangularization rigidity). This is a purely finite-dimensional matrix-algebra functorial argument, involving no $N\to\infty$ limit or analytic estimate.
\item \textbf{(3) $\Rightarrow$ (4)}: By the construction of the spectral lift operator $\mathcal{L}_{1/2}^{(N)}$ in Section~5.1, $\mathcal{H}_{1/2}^{(N)}=0$ locks the real part at $1/2$. Combined with Theorem~\ref{thm:weilconv}, the frequency measure $\nu_{1/2}^{(N)}$ is by construction supported on the critical line, and its weak limit inherits this collinearity.
\item \textbf{(4) $\Rightarrow$ (1)}: By Proposition~\ref{prop:repulsion} (constructive repulsion of off-axis zeros). If an off-axis zero $\rho_0=\sigma_0+i\gamma_0$ existed, the constructive support of $\nu_{1/2}^{(N)}$ (strictly on the critical line) would directly contradict the exponential off-axis growth required by the explicit formula. Hence all zeros must lie on the critical line.
\item \textbf{(1) $\Rightarrow$ (2)}: If RH holds, then the self-duality of the Riemann functional equation on the critical line forces $\mathcal{U}_{1/2}=\mathbb{I}$ (Theorem~\ref{thm:selfdual}), and $\mathcal{F}_{1/2}^{(N)}$ reduces to the standard block DFT; the FFT butterfly operation directly yields the recursive block-triangularization normal form.
\end{itemize}
\end{proof}

\subsection{Strict Derivation of RH (Double-Defense-Line Argument)}

By Lemma~\ref{lem:rigid} and Theorem~\ref{thm:selfdual}, at $\alpha=1/2$ the scaled generator is strictly zero:
\[
\mathcal{H}_{1/2}^{(N)}=0.
\]
This corresponds to freezing of the internal degrees of freedom---the real-part deviation is locked to zero, and all spectral points are constrained to the line $\Re(s)=1/2$. However, this only determines the real-coordinate of the zeros. To recover the full complex zeros $\rho_j = 1/2 + i\gamma_j$ (including imaginary-part information), we must introduce the external frequency indices $t_k = \frac{2\pi k}{\log N}$, realized through the frequency spectral measure $\nu_{1/2}^{(N)}$:
\[
\nu_{1/2}^{(N)} := \frac{1}{\log N}\sum_{k=0}^{N-1} |D_N(t_k)|^2 \delta_{\frac{1}{2} + i t_k}.
\]
The support of this measure is:
\[
\mathrm{supp}(\nu_{1/2}^{(N)}) = \left\{\frac{1}{2} + i t_k \;\middle|\; k=0,\dots,N-1\right\} \subset \left\{s\in\mathbb{C} \;\middle|\; \Re(s)=\frac{1}{2}\right\}.
\]
That is, all mass is strictly distributed on the critical line $\Re(s)=1/2$. By Theorem~\ref{thm:weilconv}, as $N\to\infty$ this measure converges weakly to the empirical distribution of Riemann zeros:
\[
\lim_{N\to\infty} \nu_{1/2}^{(N)} = \frac{1}{2\pi}\sum_{\rho} \delta_{\rho}.
\]

\textbf{First defense line (constructive repulsion, finite $N$)}: Proposition~\ref{prop:repulsion} shows that if an off-axis zero $\rho_0=\sigma_0+i\gamma_0$ ($\sigma_0\neq 1/2$) existed, then the amplitude of $D_N(t)$ at $t=\gamma_0$ would grow exponentially at rate $N^{\sigma_0-1/2}$, far exceeding the $(\log N)^2$-order peaks of critical-line zeros. This would force $\nu_{1/2}^{(N)}$ to carry an exponentially exploding off-axis component at grid points in a neighborhood of $\gamma_0$, directly contradicting the strict critical-line support of $\nu_{1/2}^{(N)}$. Thus the existence of off-axis zeros is constructively excluded already at finite $N$.

\textbf{Second defense line (limit argument, $N\to\infty$)}: Even setting aside the above constructive contradiction, by lower semicontinuity of weak convergence the support of the limit measure satisfies:
\[
\operatorname{supp}\left(\sum_{\rho} \delta_{\rho}\right) \subseteq \overline{\bigcup_{N} \operatorname{supp}(\nu_{1/2}^{(N)})} \subseteq \left\{s \in \mathbb{C} \;\middle|\; \Re(s)=\frac{1}{2}\right\}.
\]
If an off-axis zero $\rho_0$ ($\Re(\rho_0)\neq 1/2$) existed, the support of $\delta_{\rho_0}$ would fall outside the critical line, contradicting the above inclusion. This limit argument does not rely on boundedness of weights at finite $N$, and independently blocks any possibility of ``limit rescue''.

Therefore, all non-trivial zeros $\rho$ must satisfy $\Re(\rho)=1/2$.
\end{proof}

\textbf{Key distinction}:
\begin{itemize}
\item \textbf{Internal spectral measure} $\mu_{1/2}^{(N)}=\delta_0$: reflects $\mathcal{H}_{1/2}^{(N)}=0$, signifying only real-part locking (single-point compactification), containing no imaginary-part information.
\item \textbf{Frequency spectral measure} $\nu_{1/2}^{(N)}$: reflects imaginary-part information unfolded through FFT frequencies $t_k$, supported on the critical line, converging to the zero distribution. The exponential growth of off-axis zeros directly conflicts with this support.
\end{itemize}
Together they constitute the complete spectral lift: the ``internal'' part of $\mathcal{L}_{1/2}^{(N)}$ (from $\mathcal{H}$) locks the real part at $1/2$, while the ``external'' part (from $\mathrm{diag}(t_k)$) unfolds the imaginary part to $\gamma_j$. Only at $\alpha=1/2$ does this dual structure achieve strict self-dual matching.

\begin{corollary}
This framework provides a constructive proof pathway: approximating infinite-dimensional analytic objects (Riemann zeros) through finite-dimensional self-dual computations (standard DFT/FFT), with all steps satisfying Bishop's constructivity requirements. The NCDFT framework reveals that the Riemann Hypothesis is essentially a duality condition: only when the parameter $\alpha=1/2$ does the non-commutative discrete Fourier transform achieve strict self-duality, where the functorial fidelity between arithmetic and spectrum reaches perfect matching, and the spectral measure compactifies from the non-compact extended state ($\alpha \neq 1/2$) to the discrete point state on the critical line ($\alpha = 1/2$). The recursive block-triangularization rigidity (Lemma~\ref{lem:rigid}) hardens this analytic duality into purely algebraic functorial rigidity---the NCDFT family generated by arithmetic data admits only the global section $\alpha=1/2$, while the constructive repulsion of Proposition~\ref{prop:repulsion} and the limit support inclusion relation form a double defense line ensuring that the zero distribution, as the spectral corollary of this section, is necessarily collinear.
\end{corollary}

\appendix
\section{FFT Butterfly Operation and Duality Structure}

\subsection{Decimation-in-Time FFT Butterfly}

Consider the length-$N$ DFT ($N=2^m$):
\[
X_k = \sum_{n=0}^{N-1} x_n \omega_N^{kn}, \quad \omega_N = e^{-2\pi i/N}.
\]
Separate the input sequence into even and odd indices:
\[
X_k = \sum_{m=0}^{N/2-1} x_{2m} \omega_N^{k(2m)} + \sum_{m=0}^{N/2-1} x_{2m+1} \omega_N^{k(2m+1)} = E_k + \omega_N^k O_k,
\]
where $E_k$ and $O_k$ are length-$N/2$ DFTs:
\[
E_k = \sum_{m=0}^{N/2-1} x_{2m} \omega_{N/2}^{km}, \quad O_k = \sum_{m=0}^{N/2-1} x_{2m+1} \omega_{N/2}^{km}.
\]
For $k$ and $k+N/2$, using $\omega_N^{k+N/2} = -\omega_N^k$, we obtain the classic butterfly:
\[
\begin{cases}
X_k = E_k + \omega_N^k O_k, \\
X_{k+N/2} = E_k - \omega_N^k O_k.
\end{cases}
\]
This operation can be pictured as a butterfly structure, hence the name.

\subsection{Butterfly Operation and Duality}

The butterfly operation reveals the intrinsic symmetry of the DFT: the even part $E_k$ and odd part $O_k$ are combined via the twiddle factor $\omega_N^k$, and the twiddle factor itself satisfies $\omega_N^{k+N/2} = -\omega_N^k$, which is precisely the discrete manifestation of the $s \leftrightarrow 1-s$ symmetry in the functional equation. In the NCDFT framework, the duality $\mathcal{F}_\alpha \leftrightarrow \mathcal{F}_{1-\alpha}$ corresponds to the interchange of parameters $\alpha$ and $1-\alpha$, while the even--odd decomposition and sign reversal in the butterfly operation exactly mimic this duality.

Furthermore, viewing the NCDFT matrix $\mathcal{F}_\alpha^{(N)}$ as a DFT with internal phases $\exp(i(\alpha-1/2)\operatorname{Li}(x_n)\mathbf{H})$, its fast algorithm can similarly be decomposed into analogous butterfly structures, except that the twiddle factors must be multiplied by the corresponding diagonal matrices. This structure guarantees the computational efficiency of NCDFT while preserving the dual link to the Riemann functional equation.

\subsection{Butterfly Operation and the Functional Equation}

Combining the butterfly operation with the correspondence table of Section~4.2:

\begin{center}
\begin{tabular}{|l|l|}
\hline
FFT butterfly operation & NCDFT duality \\
\hline
Even--odd decomposition $E_k, O_k$ & Parameter duality $\alpha \leftrightarrow 1-\alpha$ \\
Twiddle factors $\omega_N^k$ and $-\omega_N^k$ & Phase factors $\exp(i(\alpha-1/2)\operatorname{Li}(x_n)\mathbf{H})$ and its conjugate \\
Symmetry of $X_k$ and $X_{k+N/2}$ & Self-dual condition $\mathcal{U}_{1/2}=\mathbb{I}$ \\
Grid refinement $N\to\infty$ & Sampling completeness $t_k \to \gamma_j$ (Section~4.4) \\
\hline
\end{tabular}
\end{center}

\textbf{Connection to Lemma~\ref{lem:rigid}}: The recursion in Appendix~A.1 requires that $E_k$ and $O_k$ must be DFTs of the same type of size $N/2$; this is the algebraic root of FFT butterfly feasibility. Lemma~\ref{lem:rigid} imposes exactly this recursive requirement on the NCDFT family: when $\alpha\neq 1/2$, the phase weights $W_n$ fail $W_{2m}^{(N)}=W_m^{(N/2)}$, so the even and odd parts cannot be reduced to an NCDFT of size $N/2$, and the recursive chain breaks. Therefore, the functorial compatibility of the FFT butterfly operation uniquely locks $\alpha=1/2$.

\begin{thebibliography}{9}
\bibitem{edwards} H. M. Edwards, \textit{Riemann's Zeta Function}, Dover, 2001 [Original: 1974].
\bibitem{titchmarsh} E. C. Titchmarsh, \textit{The Theory of the Riemann Zeta-Function}, 2nd ed., Oxford University Press, 1986.
\bibitem{weil} A. Weil, ``Sur les formules explicites de la th\'eorie des nombres premiers'', Comm. Sem. Math. Univ. Lund, 1952.
\bibitem{ingham} A. E. Ingham, ``On the estimation of $N(\sigma,T)$'', Quart. J. Math., 1940.
\bibitem{cooleytukey} J. W. Cooley, J. W. Tukey, ``An algorithm for the machine calculation of complex Fourier series'', Math. Comp., 1965.
\bibitem{bishop} E. Bishop, \textit{Foundations of Constructive Analysis}, McGraw-Hill, 1967.
\end{thebibliography}

\end{document}

